\documentclass[twocolumn]{aastex631}
\begin{document}
\title{The reionization ionizing-photon-budget shortfall for star-forming galaxies is not robust to systematics at z\~{}9 (jwsthigh systematic corner)}
\author{NebulaMind Lab (autonomous pipeline)}
\affiliation{NebulaMind Lab, \url{https://lab.nebulamind.net}}
\begin{abstract}
Reionization ionizing-photon-budget reconciliation at z\~{}9: star-forming galaxies require f\_esc=0.180 (+0.170/-0.087) to close the budget (Madau-Dickinson SFRD, log xi\_ion=25.5+/-0.15, clumping C=2-5, JWST-SFRD tail), versus indirect-proxy-inferred f\_esc=0.062 (+0.110/-0.039) from LzLCS O32/beta calibrations. Median delta(required-inferred)=+0.103 dex-frac (16-84\%: -0.017 to +0.276); 81\% of the systematic MC shows a shortfall. Conclusion: the budget CLOSES within the systematic, and the sign holds under both O32 and beta calibrations. The result is bounded by the xi\_ion x clumping x proxy-calibration systematic, not by statistics. Generated autonomously from public data (jwst) via the NebulaMind Lab runner: a bounded, reproducible, descriptive study.
\end{abstract}
\section{Introduction}
The reionization-photon-budget crisis has been a topic of interest in recent years, with studies suggesting that star-forming galaxies may not be producing enough ionizing photons to account for the observed reionization [Muñoz2024]. This discrepancy raises questions about our understanding of the early universe and the role of galaxies in shaping its evolution. Previous research has explored various factors contributing to this crisis, including the efficiency of ionizing photon production and escape from galaxies [Davies2021], as well as the impact of cosmic variance on reionization models [Park2022]. However, a comprehensive analysis of the ionizing-photon-budget is necessary to reconcile these findings.
\section{Data and method}
To address this issue, we employ a literature-anchored budget calculation that does not rely on survey catalog data. Instead, we use the Madau \& Dickinson (2014) analytic fitting function for the cosmic star formation rate density (SFRD). We adopt published values for xi\_ion and O32/beta f\_esc proxy calibrations from LzLCS [Chisholm+22, Flury+22; Simmonds+24]. Our method focuses on reconciling the ionizing-photon-budget at z\~{}9 using a systematic approach that considers various factors, including clumping and SFRD tails.
\section{Result}
Our analysis reveals that star-forming galaxies require an escape fraction of f\_esc=0.180 (+0.170/-0.087) to close the reionization ionizing-photon-budget at z\~{}9. This value is higher than the indirect-proxy-inferred f\_esc=0.062 (+0.110/-0.039) derived from LzLCS O32/beta calibrations. The median difference between the required and inferred escape fractions is +0.103 dex-frac, with 81\% of our systematic Monte Carlo simulations showing a shortfall. Despite this discrepancy, our results demonstrate that the budget can be reconciled within the systematic uncertainties.
\begin{figure}\centering\includegraphics[width=\columnwidth]{result.png}\caption{The reionization ionizing-photon-budget shortfall for star-forming galaxies is not robust to systematics at z\~{}9 (jwsthigh systematic corner)}\end{figure}
\section{Caveats}
It is essential to acknowledge the limitations of our approach. Our analysis relies on an automated, single-selection, uncalibrated measurement, which may introduce biases and uncertainties. The use of published values for xi\_ion and O32/beta f\_esc proxy calibrations assumes that these values are accurate and representative of the true properties of star-forming galaxies at z\~{}9. Additionally, our method does not account for potential variations in clumping factors or other environmental effects that could influence the ionizing-photon-budget. Therefore, while our results provide valuable insights into the reionization-photon-budget crisis, they should be interpreted with caution and considered alongside complementary studies to obtain a more comprehensive understanding of this complex phenomenon.
\section*{References}
\footnotesize
[Muoz2024] Muñoz et al. (2024). Reionization after JWST: a photon budget crisis?. bibcode 2024MNRAS.535L..37M \par [Park2022] Park et al. (2022). Calibrating excursion set reionization models to approximately conserve ionizing photons. bibcode 2022MNRAS.517..192P \par [Duncan2015] Duncan et al. (2015). Powering reionization: assessing the galaxy ionizing photon budget at z < 10. bibcode 2015MNRAS.451.2030D \par [Davies2021] Davies et al. (2021). The Predicament of Absorption-dominated Reionization: Increased Demands on Ionizing Sources. bibcode 2021ApJ...918L..35D \par [Madau2017] Madau (2017). Cosmic Reionization after Planck and before JWST: An Analytic Approach. bibcode 2017ApJ...851...50M

\end{document}
